\documentclass[12pt,a4paper]{ctexart}

\usepackage{geometry}
\geometry{
    top=2.54cm,
    bottom=2.54cm,
    left=3.18cm,
    right=3.18cm
}

\usepackage{enumitem}
\setlist{noitemsep}

\usepackage{minted}
\setminted{
    fontsize=\small,
    frame=single,
    linenos=true
}

\usepackage{amsmath}

\renewcommand{\contentsname}{Contents}
\renewcommand{\listtablename}{List of Tables}
\renewcommand{\tablename}{Table}

\title{GD32F470I Development Notes}
\author{Zonghan Chen}
\date{}

\begin{document}

\maketitle

\newpage
\tableofcontents

\newpage
\listoftables

\newpage
\section{Firmware Code Reproduction}

\subsection{01 GPIO Running LED}

\begin{enumerate}
    \item Light up LED 1 and turn off LED 3 (hold for 1 second).
    \item Light up LED 2 and turn off LED 1 (hold for 1 second).
    \item Light up LED 3 and turn off LED 2 (hold for 1 second).
          
          (Loop execution)
\end{enumerate}

\textbf{GPIO modes:}
\begin{itemize}
    \item \textbf{Pull-up Input:} The internal pull-up resistor pulls the default level to high.
    \item \textbf{Pull-down Input:} The internal pull-down resistor pulls the default level to low.
    \item \textbf{Floating Input:}
          \begin{itemize}
              \item The pin level is determined by the external circuit.
              \item There is no internal pull-up/pull-down resistor.
          \end{itemize}
    \item \textbf{Analog Input:}
          \begin{itemize}
              \item The pin is directly connected to the ADC.
              \item This mode is often used for analog signal acquisition.
          \end{itemize}
    \item \textbf{Open-Drain Output:}
          \begin{itemize}
              \item The pin can only output low level or high impedance state.
              \item Require an external pull-up resistor for high level output.
          \end{itemize}
    \item \textbf{Push-Pull Output:}
          \begin{itemize}
              \item The pin can actively output high level or low level.
              \item This mode is suitable for digital signal output.
          \end{itemize}
    \item \textbf{Alternate Function Open-Drain:} Open-drain mode controlled by an on-chip peripheral.
    \item \textbf{Alternate Function Push-Pull:} Push-pull mode controlled by an on-chip peripheral.
\end{itemize}

\subsection{02 GPIO Key Polling Mode}

\begin{enumerate}
    \item Press the Tamper Key to switch the LED 2 on/off state.
\end{enumerate}

\subsection{03 EXTI Key Interrupt Mode}

\begin{enumerate}
    \item LED 2 flashes once for test.
    \item Press the Tamper Key to switch the LED 2 on/off state.
\end{enumerate}

\textbf{External interrupt handler:}
\begin{itemize}
    \item All of the external interrupt handlers are only declared but not defined.
    \item When using an external interrupt, the corresponding external interrupt handler needs to be defined manually.
\end{itemize}

\subsection{04 USART Printf}

\begin{enumerate}
    \item The data receiving area of the serial port debugging tool displays "USART printf example: please press the Tamper key".
    \item After pressing the Tamper Key, the data receiving area of the serial port debugging tool displays "USART printf example" (repeatable operation).
\end{enumerate}

\textbf{Retarget the C library printf function to the USART:}
\begin{minted}{c}
int fputc(int ch, FILE *f)
{
    usart_data_transmit(EVAL_COM0, (uint8_t)ch);

    while (usart_flag_get(EVAL_COM0, USART_FLAG_TBE) == RESET)
    {
    }

    return ch;
}
\end{minted}

\subsection{05 USART Echo Interrupt Mode}

\begin{enumerate}
    \item LED 1, LED 2 and LED 3 flash three times for test.
    \item The data receiving area of the serial port debugging tool displays the data in the tx\_buffer array.
    \item Send the data in data.txt to the rx\_buffer array through the data sending area of the serial port debugging tool.
    \item If the data in tx\_buffer array and rx\_buffer array are the same, all the LEDs will light up successively and then turn off simultaneously (loop execution).
    \item Else if the data in tx\_buffer array and rx\_buffer array are different, all the LEDs will flash three times (loop execution).
\end{enumerate}

\subsection{06 USART DMA}

\begin{enumerate}
    \item LED 1, LED 2 and LED 3 flash three times for test.
    \item The data receiving area of the serial port debugging tool displays the data in the tx\_buffer array.
    \item Send the data in data.txt to the rx\_buffer array through the data sending area of the serial port debugging tool.
    \item If the data in tx\_buffer array and rx\_buffer array are the same, all the LEDs will light up successively and then turn off simultaneously (loop execution).
    \item Else if the data in tx\_buffer array and rx\_buffer array are different, all the LEDs will flash three times (loop execution).
\end{enumerate}

\subsection{07 ADC Temperature Vrefint}

\begin{enumerate}
    \item The data receiving area of the serial port debugging tool displays the temperature, internal reference voltage and external battery voltage values every two seconds.
          
          (Loop execution)
\end{enumerate}

\textbf{Formula:}
\begin{itemize}
    \item \textbf{Temperature:}
          \begin{equation}
              Temperature(^\circ\text{C})=\frac{V_{25^\circ\text{C}}-V_{SENSE}}{AvgSlope}+25
              \label{eq:Temperature}
          \end{equation}
          \begin{table}[htbp]
              \centering
              \caption{Parameters in Equation \eqref{eq:Temperature}}
              \begin{tabular}{|c|c|}
                  \hline
                  Parameter              & Explanation                                                           \\
                  \hline
                  $V_{25^\circ\text{C}}$ & Internal temperature sensor output voltage at $25^\circ\text{C}$      \\
                  \hline
                  $V_{SENSE}$            & Internal temperature sensor output voltage at the current temperature \\
                  \hline
                  $AvgSlope$             & Average slope of the $Temperature-V_{SENSE}$ curve                    \\
                  \hline
              \end{tabular}
          \end{table}
    \item \textbf{Internal reference voltage ($V_{REFINT}$) / Temperature sensor voltage ($V_{SENSE}$) :}
          \begin{equation}
              V_{REFINT}=\frac{ADC_{value}*V_{reference}}{2^{bits}-1}
              \label{eq:Internal-reference-voltage}
          \end{equation}
          \begin{equation}
              V_{SENSE}=\frac{ADC_{value}*V_{reference}}{2^{bits}-1}
              \label{eq:Temperature-sensor-voltage}
          \end{equation}
          \begin{table}[htbp]
              \centering
              \caption{Parameters in Equation \eqref{eq:Internal-reference-voltage} and \eqref{eq:Temperature-sensor-voltage}}
              \begin{tabular}{|c|c|}
                  \hline
                  Parameter       & Explanation           \\
                  \hline
                  $ADC_{value}$   & ADC measurement value \\
                  \hline
                  $V_{reference}$ & ADC reference voltage \\
                  \hline
                  $bits$          & ADC bit number        \\
                  \hline
              \end{tabular}
          \end{table}
    \item \textbf{External battery voltage ($V_{BAT}$) :}
          \begin{equation}
              V_{BAT}=\frac{ADC_{value}*V_{reference}}{2^{bits}-1}*4
              \label{eq:External-battery-voltage}
          \end{equation}
          \begin{table}[htbp]
              \centering
              \caption{Parameters in Equation \eqref{eq:External-battery-voltage}}
              \begin{tabular}{|c|c|}
                  \hline
                  Parameter       & Explanation           \\
                  \hline
                  $ADC_{value}$   & ADC measurement value \\
                  \hline
                  $V_{reference}$ & ADC reference voltage \\
                  \hline
                  $bits$          & ADC bit number        \\
                  \hline
              \end{tabular}
          \end{table}
          \begin{itemize}
              \item There is a bridge divider by 4 integrated on the $V_{BAT}$ pin, it connects $V_{BAT}/4$ to the ADC\_IN18 input channel. Therefore, the converted digital value of ADC\_IN18 input channel is $V_{BAT}/4$.
          \end{itemize}
\end{itemize}

\subsection{08 ADC0 ADC1 Follow Up Mode}

\begin{enumerate}
    \item The data receiving area of the serial port debugging tool displays the regular value of ADC0 and ADC1 by adc\_value[0] and adc\_value[1] every two seconds.
          
          (Loop execution)
\end{enumerate}

\textbf{The working principle of the TIMER Prescaler and Auto-Reload Register:}
\begin{itemize}
    \item \textbf{Prescaler:}
          \begin{itemize}
              \item \textbf{Hardware design logic:} A prescaler is a counter that starts counting from zero to the set prescaler value, then returns to zero and generates a pulse.
              \item \textbf{Configuration method:} If the frequency division coefficient is required to be $N$, the value of \verb|timer_parameter_struct.prescaler| should be set to $N-1$.
          \end{itemize}
    \item \textbf{Auto-Reload Register (APR) :}
          \begin{itemize}
              \item \textbf{Role:}
                    \begin{itemize}
                        \item Determine the overflow period of the timer.
                        \item The timer starts counting from zero and increments by one for each received clock pulse. When the count reaches the value of ARR, the timer overflows and triggers an interrupt or updates an event, then starts counting from zero again.
                    \end{itemize}
              \item \textbf{Configuration method:} If the timer overflow period is required to be $N$, the value of \verb|timer_parameter_struct.period| should be set to $N-1$.
          \end{itemize}
\end{itemize}

\textbf{PWM output mode:}
\begin{table}[htbp]
    \centering
    \caption{PWM Output Mode}
    \begin{tabular}{|c|c|c|}
        \hline
                                             & PWM MODE 0                   & PWM MODE 1                   \\
        \hline
        Condition of output the set polarity & $Value_{CNT}<Value_{preset}$ & $Value_{CNT}>Value_{preset}$ \\
        \hline
    \end{tabular}
\end{table}

\subsection{09 ADC0 ADC1 Regular Parallel Mode}

\begin{enumerate}
    \item The data receiving area of the serial port debugging tool displays the regular value of ADC0 and ADC1 by adc\_value[0] and adc\_value[1] every two seconds.
          
          (Loop execution)
\end{enumerate}

\textbf{ADC sync mode:}
\begin{itemize}
    \item \textbf{Free mode:} Each ADC works independently and does not interfere with each other.
    \item \textbf{Routine parallel mode:} All of the ADCs convert the routine sequence parallelly at the selected external trigger of ADC0.
    \item \textbf{Routine follow-up mode:} ADC0 converts the routine sequence at the selected external trigger. After a delay time, ADC1 converts the routine sequence. After another delay time, ADC2 converts the routine sequence.
\end{itemize}

\subsection{24 TLI IPA}

\begin{enumerate}
    \item The LCD displays a running leopard with the logo of GigaDevice as the background.
\end{enumerate}

\subsection{old 01 Starting LED Master}

\begin{enumerate}
    \item LED 1 lights up for one second and turns off for one second.
    \item LED 2 lights up for one second and turns off for one second.
    \item LED 3 lights up for one second and turns off for one second.
          
          (Loop execution)
\end{enumerate}

\subsection{old 02 Keyboard Polling Mode}

\begin{enumerate}
    \item Press the Tamper Key to switch the LED 1 on/off state.
\end{enumerate}

\subsection{old 03 Key External Interrupt Mode}

\begin{enumerate}
    \item Light up LED 1.
    \item Press the Tamper Key to switch the LED 2 on/off state.
\end{enumerate}

\end{document}